% Compilar con XeLaTeX o LuaLaTeX (en Overleaf: Menú > Compilador).
\documentclass[12pt,letterpaper]{article}

\usepackage{iftex}
\ifPDFTeX
  \usepackage[utf8]{inputenc}
  \usepackage[T1]{fontenc}
  \usepackage{lmodern}
  \usepackage[spanish]{babel}
\else
  \usepackage{fontspec}
  \usepackage{polyglossia}
  \setmainlanguage{spanish}
\fi

\usepackage[margin=2.5cm]{geometry}
\usepackage{amsmath}
\usepackage{tikz}
\usetikzlibrary{arrows.meta,angles,quotes}
\usepackage[font=small]{caption}

\pagestyle{empty}

\title{Reto III.3: asistencia gravitacional en Júpiter\\ Introducción, objetivos y marco teórico}
\author{Frankly Giovanni Poveda Pinzón \and Camila Castillo López}
\date{2 de octubre de 2026}

\setlength{\emergencystretch}{3em}
\begin{document}
\maketitle
\thispagestyle{empty}

\section{Introducción}

Llevar una sonda a los planetas exteriores exige mucha energía. Una transferencia de
Hohmann directa de la Tierra a Saturno tarda unos seis años, y acortar ese tiempo con
propulsión requiere más combustible del que un lanzador puede poner en órbita. La solución
que se usa desde los años setenta es el sobrevuelo planetario o asistencia gravitacional:
la sonda pasa cerca de un planeta y sale con una velocidad distinta, en magnitud y
dirección, respecto al Sol, sin gastar combustible. Las sondas Voyager aprovecharon en 1977
una alineación de Júpiter y Saturno para llegar a Saturno en cuatro años o menos
\cite{kohlhase1977}.

El fenómeno parece violar la conservación de la energía, pues la sonda sale de Júpiter con
más rapidez de la que traía. La explicación está en el sistema de referencia: respecto a
Júpiter la rapidez de llegada y de salida son iguales, pero Júpiter se mueve respecto al Sol,
y al sumar su velocidad la rapidez heliocéntrica de la sonda sí cambia. Entender qué se
conserva y dónde es justamente el tema de este trabajo.

El problema completo es de varios cuerpos y no tiene solución analítica. En el diseño de
misiones se usa la aproximación de cónicas empalmadas, que divide la trayectoria en tramos
donde actúa un solo cuerpo y en cada uno resuelve un problema de Kepler. La aproximación es
útil porque las escalas del problema están muy separadas, pero no es exacta. En este trabajo
se diseñará un sobrevuelo de Júpiter que lleve una sonda hasta Saturno, se resolverá con
cónicas empalmadas y con integración numérica de $N$ cuerpos, y se comparará el resultado de
ambos métodos para saber en qué condiciones dejan de coincidir.

\section{Objetivos}

\subsection*{Objetivo general}

Diseñar una asistencia gravitacional en Júpiter que lleve una sonda hasta Saturno y
determinar, comparando la aproximación de cónicas empalmadas con la integración de $N$
cuerpos, qué cantidades se conservan en cada sistema de referencia y cuál es el error de la
aproximación.

\subsection*{Objetivos específicos}

\begin{enumerate}
\item Plantear el problema en los sistemas de referencia heliocéntrico, jovicéntrico y
rotante Sol--Júpiter, e identificar las cantidades conservadas en cada uno.
\item Construir un modelo de cónicas empalmadas (problema de Lambert más hipérbola de
sobrevuelo) y encontrar con él una trayectoria Tierra--Júpiter--Saturno sin propulsión
durante el encuentro.
\item Integrar numéricamente el mismo encuentro con un código de $N$ cuerpos y verificar la
conservación de la energía y de la constante de Jacobi.
\item Cuantificar la diferencia entre ambos modelos en función de la velocidad de llegada y
del radio de empalme.
\item Comparar los resultados con la trayectoria real de Voyager~2.
\end{enumerate}

\section{Marco teórico}

\subsection{Problema de dos cuerpos}

Dos masas que interactúan gravitacionalmente se reducen a un solo cuerpo de masa reducida
que se mueve en el potencial $-\mu/r$, con $\mu = G(m_1+m_2)$. Como la fuerza es central se
conservan el momento angular específico $\vec h = \vec r\times\vec v$, lo que confina el
movimiento a un plano, y la energía específica
\begin{equation}
\varepsilon = \frac{v^2}{2} - \frac{\mu}{r}.
\end{equation}
La trayectoria es una cónica con el centro de fuerza en uno de sus focos,
\begin{equation}
r(f) = \frac{h^2/\mu}{1 + e\cos f},
\end{equation}
donde $f$ es la anomalía verdadera y $e$ la excentricidad. El signo de la energía define el
tipo de órbita: elipse si $\varepsilon<0$, parábola si $\varepsilon=0$ e hipérbola si
$\varepsilon>0$. Con el semieje mayor $a = -\mu/2\varepsilon$ la conservación de la energía
toma la forma de la ecuación de vis-viva,
\begin{equation}
v^2 = \mu\left(\frac{2}{r} - \frac{1}{a}\right),
\end{equation}
válida para las tres cónicas si en la hipérbola se toma $a<0$ \cite{curtis2020}.

\subsection{Transferencias y problema de Lambert}

La transferencia más económica entre dos órbitas circulares coplanares es la de Hohmann: una
semielipse tangente a ambas. Entre la Tierra y Júpiter su semieje es $a_t = (r_\oplus +
r_J)/2$ y el viaje dura medio periodo, unos 2.7 años. Para un diseño real se necesita algo
más general, porque los planetas deben estar en el lugar correcto en el momento de llegada.
Ese es el problema de Lambert: dadas dos posiciones $\vec r_1$ y $\vec r_2$ y un tiempo de
vuelo $\Delta t$, encontrar la órbita que las une. El teorema de Lambert establece que
$\Delta t$ depende solo del semieje mayor, de la suma $r_1+r_2$ y de la cuerda
$c=|\vec r_2-\vec r_1|$, y el problema se resuelve numéricamente, por ejemplo con la
formulación de variable universal \cite{curtis2020}.

\subsection{La hipérbola de sobrevuelo}

Cerca de Júpiter, y si se ignora al Sol, la sonda está en un problema de dos cuerpos con
$\varepsilon>0$, así que describe una hipérbola. Lejos del planeta su rapidez tiende a un
valor $v_\infty$, con $\varepsilon = v_\infty^2/2$. Como la energía se conserva, la rapidez
$v_\infty$ es la misma a la entrada y a la salida; el encuentro solo cambia la dirección
de la velocidad.

Las asíntotas corresponden a $1+e\cos f_\infty=0$, es decir $\cos f_\infty = -1/e$. El
ángulo que gira la velocidad entre la entrada y la salida es $\delta = 2f_\infty - \pi$, de
donde
\begin{equation}
\sin\frac{\delta}{2} = \frac{1}{e}.
\end{equation}
Evaluando la vis-viva en el periapsis $r_p$ se obtiene la excentricidad en términos de las
cantidades que se pueden escoger en el diseño,
\begin{equation}
e = 1 + \frac{r_p v_\infty^2}{\mu_J}.
\end{equation}
Un paso más cercano ($r_p$ pequeño) o más lento ($v_\infty$ pequeño) produce un giro mayor.
El cambio vectorial de velocidad tiene magnitud $|\Delta\vec v| = 2v_\infty\sin(\delta/2)$.
La figura~\ref{fig:esquema}a muestra esta geometría.

\begin{figure}[h]
\centering
\begin{tikzpicture}[scale=0.62, >=Stealth]
  % (a) hipérbola con foco en Júpiter, e = 1.5, r_p = 1
  \draw[thick] plot[domain=-112:112, samples=120, smooth]
     ({2.5/(1+1.5*cos(\x))*cos(\x)}, {2.5/(1+1.5*cos(\x))*sin(\x)});
  \fill (0,0) circle (3pt) node[below left] {J};
  \draw[dashed] (0,0) -- node[below] {$r_p$} (1,0);
  % asíntotas por el centro (3,0)
  \draw[gray, thin] (3,0) -- ++({-0.667*7.5},{0.745*7.5});
  \draw[gray, thin] (3,0) -- ++({-0.667*7.5},{-0.745*7.5});
  \draw[gray, thin, dotted] (3,0) -- ++({0.667*2.4},{0.745*2.4});
  \coordinate (C) at (3,0);
  \coordinate (A) at ({3+0.667},{0.745});
  \coordinate (B) at ({3-0.667},{0.745});
  \pic[draw, "$\delta$", angle radius=0.9cm, angle eccentricity=1.4] {angle = A--C--B};
  % sentido de movimiento
  \draw[->, thick] (-1.0,-4.2) -- ++(0.32,0.36) node[right] {$\vec v_{\infty,\mathrm{ent}}$};
  \draw[->, thick] (-1.2,4.4) -- ++(-0.4,0.45) node[left] {$\vec v_{\infty,\mathrm{sal}}$};
  \node at (-3.5,5.6) {(a)};

  % (b) triángulo de velocidades
  \begin{scope}[xshift=7.5cm, yshift=1cm]
    \coordinate (O) at (0,0);
    \coordinate (J) at (5,0);
    \coordinate (I) at ({5+2.6*cos(-130)},{2.6*sin(-130)});
    \coordinate (S) at ({5+2.6*cos(-46.4)},{2.6*sin(-46.4)});
    \draw[gray, dashed] (J) circle (2.6);
    \draw[->, very thick] (O) -- node[above] {$\vec v_J$} (J);
    \draw[->] (J) -- node[left, xshift=-2pt] {$\vec v_{\infty,\mathrm{ent}}$} (I);
    \draw[->] (J) -- node[right] {$\vec v_{\infty,\mathrm{sal}}$} (S);
    \draw[->, thick, dashed] (O) -- node[below left] {$\vec v_{\mathrm{antes}}$} (I);
    \draw[->, thick] (O) -- node[below, sloped] {$\vec v_{\mathrm{después}}$} (S);
    \fill (O) circle (2pt) node[left] {Sol};
    \node at (0,4.6) {(b)};
  \end{scope}
\end{tikzpicture}
\caption{(a) Trayectoria hiperbólica alrededor de Júpiter (J) en su propio sistema de
referencia. Las líneas grises son las asíntotas y $\delta$ el ángulo que gira la velocidad.
(b) Suma de velocidades vista desde el Sol. Como $|\vec v_\infty|$ no cambia, la velocidad de
salida termina sobre el círculo punteado; si $\vec v_\infty$ gira hacia $\vec v_J$, la
rapidez heliocéntrica aumenta.}
\label{fig:esquema}
\end{figure}

\subsection{Cambio de sistema de referencia}

En el sistema del Sol la velocidad de la sonda es $\vec v = \vec v_J + \vec v_\infty$. Si
durante el encuentro la velocidad de Júpiter no cambia apreciablemente, el cambio en la
energía heliocéntrica de la sonda es
\begin{equation}
\Delta\varepsilon_\odot = \tfrac12|\vec v_J+\vec v_{\infty,\mathrm{sal}}|^2
                       - \tfrac12|\vec v_J+\vec v_{\infty,\mathrm{ent}}|^2
                       = \vec v_J\cdot\left(\vec v_{\infty,\mathrm{sal}}
                         - \vec v_{\infty,\mathrm{ent}}\right),
\end{equation}
donde se usó que $|\vec v_{\infty,\mathrm{sal}}| = |\vec v_{\infty,\mathrm{ent}}|$. La sonda
gana energía cuando el giro acerca $\vec v_\infty$ a la dirección de movimiento de Júpiter,
lo que ocurre si pasa por detrás del planeta, y la pierde en el caso contrario. La energía
ganada sale de la órbita de Júpiter; como la razón de masas es del orden de $10^{-25}$, el
cambio en Júpiter es imperceptible. En el sistema heliocéntrico, entonces, ni la energía
ni el momento angular de la sonda se conservan, aunque sí los del sistema completo.

\subsection{Esfera de influencia y cónicas empalmadas}

Para decidir dónde cambiar del sistema del Sol al de Júpiter se usa la esfera de influencia.
Laplace la definió comparando, en cada sistema de referencia, la perturbación del tercer
cuerpo con la atracción del cuerpo principal: dentro de la esfera la perturbación del Sol
sobre el movimiento relativo a Júpiter es menor, en proporción, que la de Júpiter sobre el
movimiento relativo al Sol. El resultado es
\begin{equation}
r_{\mathrm{SOI}} = a_J\left(\frac{m_J}{M_\odot}\right)^{2/5},
\end{equation}
que para Júpiter vale unos $4.8\times10^7$~km \cite{murray1999}. El método de cónicas
empalmadas resuelve un problema de Kepler heliocéntrico fuera de la esfera y uno jovicéntrico
dentro de ella, igualando posición y velocidad en la frontera. Visto desde el Sol, como la
esfera es pequeña frente a la órbita, el encuentro se puede tratar como instantáneo.

El criterio no dice que la atracción del Sol sea pequeña dentro de la esfera, sino que su
efecto diferencial, la diferencia entre lo que el Sol acelera a la sonda y a Júpiter, lo es.
Por eso la aproximación debería empeorar cuando la sonda pasa mucho tiempo cerca de la
frontera.

\subsection{Problema restringido de tres cuerpos}

Si Júpiter describe un círculo alrededor del centro de masa y la sonda no perturba a ninguno
de los dos cuerpos, se tiene el problema restringido circular de tres cuerpos. En un sistema
que rota con velocidad angular $n$ junto con Júpiter, las ecuaciones de movimiento de la sonda
son
\begin{equation}
\ddot x - 2n\dot y = \frac{\partial U}{\partial x}, \qquad
\ddot y + 2n\dot x = \frac{\partial U}{\partial y}, \qquad
U = \frac{n^2}{2}(x^2+y^2) + \frac{\mu_1}{r_1} + \frac{\mu_2}{r_2},
\end{equation}
con $r_1$ y $r_2$ las distancias al Sol y a Júpiter. El potencial no depende del tiempo en
este sistema, y multiplicando por $\dot x$ y $\dot y$ y sumando se obtiene una integral de
movimiento, la constante de Jacobi
\begin{equation}
C_J = 2U - (\dot x^2 + \dot y^2).
\end{equation}
Lejos de Júpiter, $C_J$ se puede escribir en términos de los elementos orbitales
heliocéntricos de la sonda, lo que da el parámetro de Tisserand
\begin{equation}
T_J = \frac{a_J}{a} + 2\cos i\,\sqrt{\frac{a}{a_J}\left(1-e^2\right)} ,
\end{equation}
que debe ser aproximadamente el mismo antes y después del encuentro \cite{murray1999}. Esta
relación limita qué órbitas se pueden alcanzar con un sobrevuelo y sirve como prueba de
consistencia de cualquier cálculo.

\subsection{Integración numérica de $N$ cuerpos}

Sin las aproximaciones anteriores, la aceleración de cada cuerpo es la suma de las
atracciones de todos los demás,
\begin{equation}
\ddot{\vec r}_i = -\sum_{j\neq i} G m_j\,\frac{\vec r_i - \vec r_j}{|\vec r_i - \vec r_j|^3},
\end{equation}
y las ecuaciones se integran numéricamente. La dificultad es que durante el encuentro las
escalas de tiempo cambian en varios órdenes de magnitud: horas cerca del periapsis y años en
el crucero. Se necesitan integradores de paso adaptativo, como IAS15, incluido en el código
REBOUND, que ajusta el paso para mantener el error por debajo de la precisión de máquina
\cite{rein2012,rein2015}. Como no existe solución exacta para comparar, la calidad de la
integración se juzga por la conservación de la energía total, del momento angular total y,
en el problema restringido, de la constante de Jacobi.

\begin{thebibliography}{9}
\bibitem{kohlhase1977} C.~E. Kohlhase y P.~A. Penzo. Voyager mission description.
\emph{Space Science Reviews} 21, 77--101 (1977).
\bibitem{curtis2020} H.~D. Curtis. \emph{Orbital Mechanics for Engineering Students}.
Butterworth-Heinemann, 4.\textsuperscript{a} ed., 2020.
\bibitem{murray1999} C.~D. Murray y S.~F. Dermott. \emph{Solar System Dynamics}.
Cambridge University Press, 1999.
\bibitem{rein2012} H.~Rein y S.-F. Liu. REBOUND: an open-source multi-purpose $N$-body code
for collisional dynamics. \emph{Astronomy \& Astrophysics} 537, A128 (2012).
\bibitem{rein2015} H.~Rein y D.~S. Spiegel. IAS15: a fast, adaptive, high-order integrator
for gravitational dynamics. \emph{Monthly Notices of the Royal Astronomical Society} 446,
1424--1437 (2015).
\end{thebibliography}

\end{document}
